\documentclass[11pt]{article}
\usepackage[utf8]{inputenc}
\usepackage[T1]{fontenc}
\usepackage[spanish]{babel}
\usepackage{amsmath,amssymb}
\usepackage{graphicx}
\usepackage{geometry}
\usepackage{xcolor}
\usepackage{listings}
\usepackage{hyperref}
\usepackage{enumitem}
\usepackage{float}
\geometry{margin=1in}

\definecolor{codebg}{RGB}{245,245,245}
\lstset{
  basicstyle=\ttfamily\small,
  backgroundcolor=\color{codebg},
  breaklines=true,
  frame=single,
  columns=fullflexible,
  showstringspaces=false,
  captionpos=b
}

\title{Proyecto: Teoría de Conjuntos y Funciones \\ \large Programa en Python para operar conjuntos y verificar funciones}
\author{
Miguel Angel David Carranza Osorio \\ Carnet 24458 \\[4pt]
David Díaz \\
Ivan Morataya \\
Jose Fabian \\
David Perez
}
\date{\today}

\begin{document}
\maketitle

\section{Especificación del problema}

El proyecto consiste en aplicar los temas de teoría de conjuntos y funciones vistos en
clase mediante un programa en Python que permita:

\begin{itemize}
  \item Construir conjuntos definidos por extensión (incluyendo pares ordenados).
  \item Definir un conjunto referencial (universo) para poder calcular complementos.
  \item Operar los conjuntos con las operaciones usuales:
  \begin{itemize}
    \item Unión ($\cup$)
    \item Intersección ($\cap$)
    \item Diferencia ($\setminus$)
    \item Complemento ($\overline{A}$)
    \item Producto cartesiano ($\times$)
  \end{itemize}
  \item Agrupar subexpresiones mediante símbolos de agrupación (paréntesis).
  \item Dado el resultado de un producto cartesiano (o cualquier conjunto de pares
  ordenados), determinar si dicho conjunto representa una función, incluyendo si es
  una función \emph{total} respecto a un conjunto dominio dado.
\end{itemize}

El programa acepta las expresiones tanto en una notación simple como en notación
LaTeX, tal como se explica en la Sección \ref{sec:sintaxis}.

\section{Diseño del programa}

\subsection{Estructura general}

El programa (\texttt{conjuntos.py}) es interactivo: muestra un menú en la consola con
las siguientes opciones:

\begin{enumerate}
  \item Definir un conjunto.
  \item Definir el conjunto referencial (universo).
  \item Evaluar una operación / expresión.
  \item Listar los conjuntos definidos.
  \item Salir.
\end{enumerate}

Internamente, cada conjunto se guarda en un diccionario de Python
(\texttt{nombre~$\rightarrow$~frozenset}), donde los elementos pueden ser letras,
números o pares ordenados (tuplas de Python), lo que permite representar tanto
conjuntos simples como relaciones (conjuntos de pares).

\subsection{Sintaxis aceptada}
\label{sec:sintaxis}

Para escribir una operación sin usar comandos LaTeX se utilizan los siguientes
símbolos:

\begin{center}
\begin{tabular}{|c|l|l|}
\hline
\textbf{Símbolo} & \textbf{Operación} & \textbf{Ejemplo} \\
\hline
\texttt{+} & Unión & \texttt{A+B} \\
\texttt{*} & Intersección & \texttt{A*B} \\
\texttt{\textbackslash} & Diferencia & \texttt{A\textbackslash B} \\
\texttt{'} & Complemento (respecto al referencial) & \texttt{A'} \\
\texttt{x} & Producto cartesiano & \texttt{AxB} \\
\texttt{( )} & Agrupación & \texttt{(A+B)*C} \\
\hline
\end{tabular}
\end{center}

También se acepta la notación LaTeX de forma literal, anteponiendo (opcionalmente) la
marca \texttt{LaTeX:}. El programa traduce automáticamente:
\texttt{\textbackslash cup}~$\to$~\texttt{+},
\texttt{\textbackslash cap}~$\to$~\texttt{*},
\texttt{\textbackslash times}~$\to$~\texttt{x},
\texttt{\textbackslash setminus}~$\to$~\texttt{\textbackslash}, y
\texttt{\textbackslash overline\{X\}}~$\to$~\texttt{(X)'}.

Para verificar si un conjunto es función se usa \texttt{fun(EXPR)}, donde
\texttt{EXPR} puede ser el nombre de un conjunto ya definido (por ejemplo un conjunto
de pares como \texttt{E}) o cualquier expresión que produzca un conjunto de pares
(por ejemplo un producto cartesiano \texttt{AxB}).

\subsection{Precedencia de operadores}

Como no todas las expresiones llevan paréntesis explícitos, el programa usa la
siguiente precedencia (de mayor a menor):

\begin{enumerate}
  \item Complemento \quad $A'$
  \item Producto cartesiano \quad $A \times B$
  \item Intersección \quad $A \cap B$
  \item Unión y diferencia (misma precedencia, de izquierda a derecha) \quad
  $A \cup B$, \; $A \setminus B$
\end{enumerate}

Los paréntesis siempre pueden usarse para forzar un orden distinto.

\subsection{Verificación de función}

Dado un conjunto de pares ordenados $R$, el programa verifica primero que $R$ esté
\emph{bien definido} como relación funcional, es decir, que ningún elemento del
dominio se relacione con dos elementos distintos del codominio. Si el usuario
proporciona además el nombre de un conjunto que se usará como dominio, el programa
verifica también que $R$ sea una función \textbf{total}: que todos los elementos de
ese conjunto dominio aparezcan como primera componente de algún par en $R$. Si falta
algún elemento, se reporta cuáles son y se concluye que no es función total (aunque
sí puede estar bien definida).

\section{Ejecución del programa}

A continuación se muestran capturas de la ejecución del programa, usando los
conjuntos de ejemplo:

\[
\begin{aligned}
U &= \{a,b,c,d,e,f,g,h,i,j,k,1,2,3,4,5\} \\
A &= \{a,1,3,d,g,h,4,5\} \\
B &= \{2,1,4,e,f,g,k\} \\
C &= \{b,d,f,h,k,2,4\} \\
D &= \{\} \\
E &= \{(1,a),(2,b),(3,c)\}
\end{aligned}
\]

\subsection{Definición de los conjuntos}

En la Figura \ref{fig:cap1} se define cada conjunto usando la opción 1 del menú, y el
conjunto referencial $U$ usando la opción 2.

\begin{figure}[H]
  \centering
  \includegraphics[width=0.95\textwidth]{capturas/captura1_definicion.png}
  \caption{Definición de los conjuntos $A$, $B$, $C$, $D$, $E$ y del referencial $U$.}
  \label{fig:cap1}
\end{figure}

\begin{figure}[H]
  \centering
  \includegraphics[width=0.95\textwidth]{capturas/captura2_listado.png}
  \caption{Listado de todos los conjuntos definidos (opción 4 del menú).}
  \label{fig:cap2}
\end{figure}

\subsection{Operaciones básicas: unión, intersección, diferencia y complemento}

\begin{figure}[H]
  \centering
  \includegraphics[width=0.95\textwidth]{capturas/captura3_operaciones_basicas.png}
  \caption{Unión $A \cup C$, intersección $A \cap C$, diferencia $A \setminus B$ y
  complemento $\overline{A}$.}
  \label{fig:cap3}
\end{figure}

Verificación manual:
\begin{itemize}
  \item $A \cup C = \{1,2,3,4,5,a,b,d,f,g,h,k\}$ — correcto, ya que une todos los
  elementos de $A$ y $C$ sin repetir.
  \item $A \cap C = \{4,d,h\}$ — son los únicos elementos que $A$ y $C$ tienen en
  común.
  \item $A \setminus B = \{3,5,a,d,h\}$ — son los elementos de $A$ que no están en
  $B$.
  \item $\overline{A} = U \setminus A = \{2,b,c,e,f,i,j,k\}$ — el resto de elementos
  del universo que no pertenecen a $A$.
\end{itemize}

\subsection{Agrupación con paréntesis}

\begin{figure}[H]
  \centering
  \includegraphics[width=0.8\textwidth]{capturas/captura4_agrupacion.png}
  \caption{Operación con símbolos de agrupación: $(A \cup B) \cap C$.}
  \label{fig:cap4}
\end{figure}

$A \cup B = \{1,2,3,4,5,a,d,e,f,g,h,k\}$, y al intersecar con
$C=\{2,4,b,d,f,h,k\}$ se obtiene $\{2,4,d,f,h,k\}$, que coincide con el resultado
mostrado por el programa.

\subsection{Verificación de función: \texttt{fun(E)}}

\begin{figure}[H]
  \centering
  \includegraphics[width=0.95\textwidth]{capturas/captura5_fun_E.png}
  \caption{Verificación de si $E=\{(1,a),(2,b),(3,c)\}$ es función, usando $A$ como
  dominio de referencia.}
  \label{fig:cap5}
\end{figure}

$E$ está bien definida (ningún elemento del dominio se repite con una imagen
distinta), pero al pedir que sea función total respecto a $A$, el programa detecta
correctamente que $A$ tiene elementos (\{4,5,a,d,g,h\}) que no aparecen como primera
componente en $E$; por lo tanto $E$ no es función total sobre $A$ (sí lo sería sobre
el conjunto $\{1,2,3\}$).

\subsection{Casos con el conjunto vacío}

\begin{figure}[H]
  \centering
  \includegraphics[width=0.95\textwidth]{capturas/captura6_conjunto_vacio.png}
  \caption{Complemento de $D$ (el conjunto vacío) y diferencia $D \setminus U$.}
  \label{fig:cap6}
\end{figure}

Como $D = \emptyset$, su complemento respecto a $U$ es el mismo $U$
($\overline{\emptyset} = U$), y $D \setminus U = \emptyset$ porque no hay ningún
elemento en $D$ del cual partir.

\subsection{Producto cartesiano y verificación de función: \texttt{fun(A\texttimes B)}}

\begin{figure}[H]
  \centering
  \includegraphics[width=0.95\textwidth]{capturas/captura7_fun_AxB.png}
  \caption{Producto cartesiano $A \times B$ y verificación de si es función.}
  \label{fig:cap7}
\end{figure}

El producto cartesiano completo $A \times B$ tiene $|A|\times|B| = 8 \times 7 = 56$
pares ordenados. Como era de esperarse, \textbf{no} es una función, porque cada
elemento de $A$ aparece emparejado con los 7 elementos de $B$ (es decir, cada
elemento del dominio se relaciona con más de un elemento distinto del codominio). El
programa detecta e informa correctamente estos conflictos. Este resultado ilustra un
punto importante de la teoría: el producto cartesiano completo $A\times B$ es en
general \emph{la relación más grande posible} entre $A$ y $B$, y solo sería función
si $|B|\le 1$, ya que una función es, por definición, un subconjunto particular
(no necesariamente todo) de $A \times B$.

\section{Conclusiones}

\begin{itemize}
  \item El programa permite construir conjuntos arbitrarios (incluyendo conjuntos de
  pares ordenados) y operar sobre ellos con unión, intersección, diferencia,
  complemento y producto cartesiano, respetando la precedencia usual y permitiendo
  agrupar con paréntesis.
  \item Aceptar tanto notación simple como notación LaTeX facilita escribir las
  operaciones tal como se plantean en clase o en la guía del proyecto.
  \item La verificación de función distingue correctamente entre una relación
  ``bien definida'' (sin ambigüedades) y una función ``total'' (que cubre todo un
  dominio dado), lo cual se comprobó tanto con un conjunto de pares definido
  manualmente ($E$) como con el resultado de un producto cartesiano completo
  ($A\times B$), donde el programa identificó correctamente que este último no es
  una función.
\end{itemize}

\section{Actualización: uso de librerías y entrada por archivo}

A partir de la versión inicial, el programa se simplificó usando la librería
estándar \texttt{itertools} (función \texttt{itertools.product}) para calcular el
producto cartesiano, en lugar de escribirlo a mano con dos ciclos anidados. Esto no
está prohibido por el enunciado del proyecto, que no restringe el uso de librerías.

Además, el programa ahora puede recibir un archivo \texttt{.txt} con las
definiciones y operaciones (en el mismo formato del ejemplo de la guía), ya sea como
argumento de línea de comandos:

\begin{lstlisting}
python3 conjuntos.py ejemplo.txt
\end{lstlisting}

o cargándolo desde el menú (opción 4). El conjunto llamado \texttt{U} en el archivo
se toma automáticamente como conjunto referencial. En el caso de \texttt{fun(...)},
el archivo puede indicar el dominio directamente como segundo argumento, por ejemplo
\texttt{fun(E,A)}, sin necesidad de que el programa pregunte de forma interactiva.
Después de procesar el archivo, el programa sigue funcionando de forma interactiva
por si se quieren hacer más operaciones.

\begin{figure}[H]
  \centering
  \includegraphics[width=0.95\textwidth]{capturas/captura8_modo_archivo.png}
  \caption{Ejecución cargando todos los conjuntos y operaciones desde
  \texttt{ejemplo.txt} en un solo paso.}
  \label{fig:cap8}
\end{figure}

Los resultados obtenidos en este modo son idénticos a los mostrados en la
Sección~3, ya que se trata de las mismas operaciones sobre los mismos conjuntos.

\end{document}
